% Problems, Chapter 10: Green's Functions
% Arfken, Weber & Harris, Mathematical Methods for Physicists, 7th ed.
% Requires: amsmath, amssymb, tikz

\section*{Problems}

\Needspace{9\baselineskip}
\subsection*{10.1 One-Dimensional Problems}

\Needspace{5\baselineskip}
\par\bigskip
\noindent\textbf{10.1.1}\quad Show that
\[
G(x,t)=\begin{cases}
x, & 0\le x<t,\\
t, & t<x\le 1,
\end{cases}
\]
is the Green's function for the operator $L=-d^{2}/dx^{2}$ and the boundary
conditions $y(0)=0$, $y'(1)=0$.
\solution
For fixed $0<t<1$, write the proposed kernel as $G(x,t)=\min(x,t)$.
It is continuous at $x=t$, and its two derivatives there are
\[
G_x(t^-,t)=1,\qquad G_x(t^+,t)=0.
\]
Thus $[G_x]_{t^-}^{t^+}=-1$. A continuous, piecewise smooth function
has a delta function in its second derivative with coefficient equal to
this derivative jump. Consequently
\[
-G_{xx}=\delta(x-t).
\]
Away from $t$, both branches are linear, so there is no ordinary part
of $-G_{xx}$. The endpoint conditions also hold:
$G(0,t)=0$ and $G_x(1,t)=0$.

To check the resulting solution directly, let $f$ be continuous and split
the integral at $t=x$:
\[
y(x)=\int_0^1G(x,t)f(t)\,dt
=\int_0^x t f(t)\,dt+x\int_x^1f(t)\,dt.
\]
Differentiation gives
\begin{align*}
y'(x)&=xf(x)+\int_x^1f(t)\,dt-xf(x)
      =\int_x^1f(t)\,dt,\\
y''(x)&=-f(x).
\end{align*}
Hence $-y''=f$, $y(0)=0$, and $y'(1)=0$. Finally, the difference of
two solutions would be $ax+b$; its endpoint conditions force $b=0$
and $a=0$. Thus this kernel gives the unique solution.

\Needspace{5\baselineskip}
\par\bigskip
\noindent\textbf{10.1.2}\quad Find the Green's function for
\begin{enumerate}
\item[(a)] $Ly(x)=\dfrac{d^{2}y(x)}{dx^{2}}+y(x)$,\quad
  $y(0)=0$,\quad$y'(1)=0$.
\item[(b)] $Ly(x)=\dfrac{d^{2}y(x)}{dx^{2}}-y(x)$,\quad$y(x)$ finite for
  $-\infty<x<\infty$, \erratum{meaning bounded as $x\to\pm\infty$}.
\end{enumerate}
\solution
(a) For $y''+y=0$, choose a left solution and a right solution
\[
u(x)=\sin x,\qquad v(x)=\cos(1-x).
\]
They satisfy $u(0)=0$ and $v'(1)=0$, respectively. Continuity at $x=t$
is built into the form
\[
G(x,t)=C\begin{cases}u(x)v(t),&x<t,\\u(t)v(x),&x>t.\end{cases}
\]
The derivative jump must be $+1$ because the coefficient of $y''$ is $1$.
The needed Wronskian is
\begin{align*}
u(t)v'(t)-u'(t)v(t)
&=\sin t\sin(1-t)-\cos t\cos(1-t)\\
&=-\cos\bigl(t+(1-t)\bigr)=-\cos1.
\end{align*}
Therefore $C(-\cos1)=1$, and
\[
\boxed{G(x,t)=-\frac1{\cos1}
\begin{cases}
\sin x\cos(1-t),&0\le x\le t,\\
\sin t\cos(1-x),&t\le x\le1.
\end{cases}}
\]
Both branches solve $G_{xx}+G=0$ away from the source, are continuous
at the source, have derivative jump $1$, and satisfy the specified
endpoint conditions. These facts prove $(\partial_x^2+1)G=\delta(x-t)$.
There is no zero mode: $h(0)=0$ gives $h=A\sin x$, and $h'(1)=A\cos1=0$
gives $A=0$ since $\cos1\ne0$.

(b) \erratum{``Finite for $-\infty<x<\infty$'' must here mean bounded
as $x\to\pm\infty$. Finiteness only at each finite point would allow
both $e^x$ and $e^{-x}$ and would not select a unique Green's function.}
The bounded homogeneous solution on the left is proportional to $e^x$;
on the right it is proportional to $e^{-x}$. Impose continuity by writing
\[
G(x,t)=C\begin{cases}e^{x-t},&x<t,\\e^{t-x},&x>t.\end{cases}
\]
Its derivative jump is $-C-C=-2C$, so $-2C=1$ gives
\[
\boxed{G(x,t)=-\frac12e^{-|x-t|}.}
\]
It decays at both infinities, solves $G_{xx}-G=0$ for $x\ne t$, and
has derivative jump $1$. Thus $y(x)=\int_{\mathbb R}G(x,t)f(t)\,dt$
solves $y''-y=f$ whenever the integral is defined, for example for
bounded continuous $f$. In that example
$|y(x)|\le\|f\|_\infty\int_{\mathbb R}e^{-|x-t|}/2\,dt
=\|f\|_\infty$.
No nonzero combination $Ae^x+Be^{-x}$ is bounded at both infinities,
which proves uniqueness in the bounded class.

\Needspace{5\baselineskip}
\par\bigskip
\noindent\textbf{10.1.3}\quad Show that the function $y(x)$ defined by
Eq.~(10.26) satisfies the initial-value problem defined by Eq.~(10.24) and its
initial conditions $y(0)=y'(0)=0$.
\solution
The referenced formulas, Eqs.~(10.24) and (10.26) on textbook
pp.~453--454, are
\[
y''+y=f(x),\qquad y(x)=\int_0^x\sin(x-t)f(t)\,dt.
\]
For continuous $f$, the Leibniz rule includes both an endpoint term and
a derivative of the integrand. Applying it once gives
\begin{align*}
y'(x)&=\sin(x-x)f(x)+\int_0^x\cos(x-t)f(t)\,dt\\
&=\int_0^x\cos(x-t)f(t)\,dt.
\end{align*}
Apply the rule a second time:
\begin{align*}
y''(x)&=\cos(x-x)f(x)-\int_0^x\sin(x-t)f(t)\,dt\\
&=f(x)-y(x).
\end{align*}
This proves $y''+y=f$. At $x=0$, both displayed integrals have zero
length, so $y(0)=y'(0)=0$.

For forward evolution $x\ge0$ with source times $t\ge0$, the equivalent
causal kernel is $G(x,t)=\Theta(x-t)\sin(x-t)$.
The displayed Volterra integral itself also solves the initial-value
problem for negative $x$ when its integral is given the usual orientation. The kernel vanishes for $x<t$, is continuous
with value zero at $x=t$, and its first derivative jumps from $0$ to
$1$. Thus $(\partial_x^2+1)G=\delta(x-t)$. The difference of any two
solutions is $A\sin x+B\cos x$; its zero initial position and velocity
force $A=B=0$, so the verified solution is unique.

\Needspace{5\baselineskip}
\par\bigskip
\noindent\textbf{10.1.4}\quad Find the Green's function for the equation
\[
-\frac{d^{2}y}{dx^{2}}-\frac{y}{4}=f(x),
\]
with boundary conditions $y(0)=y(\pi)=0$.
\par\nobreak\hfill\emph{ANS.}\enspace
$G(x,t)=\begin{cases}
2\sin(x/2)\cos(t/2), & 0\le x<t,\\
2\cos(x/2)\sin(t/2), & t<x\le\pi.
\end{cases}$
\solution
For $y''+y/4=0$, take
\[
u(x)=\sin(x/2),\qquad v(x)=\cos(x/2).
\]
Here $u(0)=0$ and $v(\pi)=0$. A continuous kernel with the required
endpoint behavior is $C u(x)v(t)$ for $x<t$ and $C u(t)v(x)$ for $x>t$.
Its derivative jump is
\begin{align*}
G_x(t^+,t)-G_x(t^-,t)
&=C\left[-\frac12\sin^2(t/2)-\frac12\cos^2(t/2)\right]\\
&=-\frac C2.
\end{align*}
Because the operator is $-\partial_x^2-1/4$, this jump must be $-1$.
Therefore $C=2$ and
\[
\boxed{G(x,t)=\begin{cases}
2\sin(x/2)\cos(t/2),&0\le x\le t,\\
2\cos(x/2)\sin(t/2),&t\le x\le\pi.
\end{cases}}
\]
The branches solve the homogeneous equation off the diagonal; continuity
and the jump prove $(-\partial_x^2-1/4)G=\delta(x-t)$.

For a direct check of the source sign, put
\[
I(x)=\int_0^x\sin(t/2)f(t)\,dt,\qquad
J(x)=\int_x^\pi\cos(t/2)f(t)\,dt.
\]
Then $y=2\cos(x/2)I+2\sin(x/2)J$. Since
$I'=\sin(x/2)f$ and $J'=-\cos(x/2)f$, the two endpoint terms cancel
in the first derivative:
\[
y'=-\sin(x/2)I+\cos(x/2)J.
\]
Differentiating again gives
\[
y''=-\tfrac12\cos(x/2)I-\tfrac12\sin(x/2)J
-\bigl[\sin^2(x/2)+\cos^2(x/2)\bigr]f=-\frac y4-f.
\]
Also $I(0)=J(\pi)=0$ gives the two Dirichlet conditions.
A homogeneous solution satisfying $y(0)=0$ is $A\sin(x/2)$, and its
value at $\pi$ is $A$, so the homogeneous boundary problem has only
the zero solution.

\Needspace{5\baselineskip}
\par\bigskip
\noindent\textbf{10.1.5}\quad Construct the Green's function for
\[
x^{2}\frac{d^{2}y}{dx^{2}}+x\frac{dy}{dx}+(k^{2}x^{2}-1)y=0,
\]
subject to the boundary conditions $y(0)=0$, $y(1)=0$.
\erratum{Use the operator exactly as printed with ordinary $dt$ source
integration. Construct its inverse for nonresonant real $k$, treat $k=0$
separately, and explain the obstruction at a positive zero of $J_1(k)$.}
\solution
\erratum{The normalization of the operator matters at this singular
endpoint. The Instructor's Manual, p.~119, first divides the printed
ODE by $x$. Its symmetric kernel is the Green's function of that divided
operator. For the operator exactly as printed, with an ordinary $dt$
source integral, the kernel has an additional factor $1/t$. An ordinary
inverse also requires a nonresonant parameter. Both conventions are
specified below.}
Write
\[
D=x^2\frac{d^2}{dx^2}+x\frac d{dx}+k^2x^2-1,
\qquad
L=\frac1xD=\frac d{dx}\left(x\frac d{dx}\right)+k^2x-\frac1x.
\]
First take $k>0$ with $J_1(k)\ne0$, and let $0<t<1$. On either side of
the source, $Dy=0$ is Bessel's equation of order one. Its independent
solutions are $J_1(kx)$ and $Y_1(kx)$. Their small-$x$ behavior is
\[
J_1(kx)\sim\frac{kx}{2},\qquad
Y_1(kx)\sim-\frac{2}{\pi kx}.
\]
Thus the left homogeneous solution obeying $y(0)=0$ is
$u(x)=J_1(kx)$. Choose the right solution
\[
v(x)=Y_1(kx)-\frac{Y_1(k)}{J_1(k)}J_1(kx),\qquad v(1)=0.
\]

To determine the normalization, define $W=uv'-u'v$, with derivatives
with respect to $x$. Dividing the homogeneous ODE by $x^2$ gives
$y''=-y'/x-(k^2-x^{-2})y$. Hence
\[
W'=uv''-u''v=-\frac1x(uv'-u'v)=-\frac Wx,
\qquad W=\frac Cx.
\]
The multiple of $J_1$ in $v$ does not change $W$. Using the displayed
small-$x$ forms of $J_1,Y_1$ gives
\[
W\sim\frac{kx}{2}\frac{2}{\pi kx^2}
      -\frac k2\left(-\frac{2}{\pi kx}\right)
=\frac{2}{\pi x}.
\]
Therefore the exact Wronskian is $W(x)=2/(\pi x)$.

For $LG_L=\delta(x-t)$, the jump condition is
$t[G_{L,x}]_{t^-}^{t^+}=1$. Set
$G_L=C u(x_<)v(x_>)$, where $x_<=\min(x,t)$ and
$x_>=\max(x,t)$. Its jump is $CW(t)$, giving
$t\,C\,2/(\pi t)=1$, or $C=\pi/2$. Consequently
\[
\boxed{G_L(x,t)=\frac\pi2J_1(kx_<)
\left[Y_1(kx_>)-\frac{Y_1(k)}{J_1(k)}J_1(kx_>)\right].}
\]
This is the symmetric kernel in the manual. Since $D=xL$,
\[
D_xG_L(x,t)=x\delta(x-t)=t\delta(x-t).
\]
Thus the answer for the printed operator is
\[
\boxed{G_D(x,t)=\frac{G_L(x,t)}t,\qquad D_xG_D(x,t)=\delta(x-t).}
\]
In particular $[G_{D,x}]=1/t^2$, exactly as required by the leading
coefficient $x^2$. The two equivalent source formulas are
\[
Dy=F\quad\Longrightarrow\quad
y(x)=\int_0^1G_D(x,t)F(t)\,dt
     =\int_0^1G_L(x,t)\frac{F(t)}t\,dt.
\]
For each interior $t$, both kernels are continuous at $x=t$, vanish at
$x=1$, and tend to zero as $x\downarrow0$.

The singular endpoint also restricts allowable sources if the integrated
solution is required to vanish at zero. Sources supported inside $(0,1)$
are sufficient; more generally $F(x)=O(x^\alpha)$ for $\alpha>0$ is
sufficient. The small-$x$ kernel estimates give terms bounded by constants
times $x^{-1}\int_0^x|F(t)|dt$ and
$x\int_x^1|F(t)|t^{-2}dt$.  Up to constant multipliers, these bounds
are respectively
\[
\frac1x\int_0^x t^\alpha\,dt=\frac{x^\alpha}{\alpha+1},
\qquad
x\int_x^1t^{\alpha-2}\,dt=
\begin{cases}
\dfrac{x-x^\alpha}{\alpha-1},&\alpha\ne1,\\
-x\log x,&\alpha=1.
\end{cases}
\]
For every $\alpha>0$ both tend to zero, including the logarithmic case.
An arbitrary nonzero constant source at zero is not covered by this
boundary-value inverse.

At a positive root $J_1(k)=0$, the nonzero function $u=J_1(kx)$
satisfies both endpoint conditions. The denominator vanishes for a
reason: the homogeneous problem has a zero mode. Green's identity gives
for an admissible solution of $Dy=F$
\[
\int_0^1u(x)\frac{F(x)}x\,dx
=\int_0^1(uLy-yLu)\,dx
=\left[x(uy'-u'y)\right]_0^1=0.
\]
A general source does not satisfy this compatibility condition, and a
compatible solution is not unique because one can add $cu$.
Therefore no ordinary inverse Green's function exists at these resonances.

The case $k=0$ must be treated separately, rather than counted as a
positive Bessel zero. Now the homogeneous solutions are $x$ and $1/x$.
Choose $u=x$, $v=x-1/x$; then $W=2/x$. The kernels are
\[
\boxed{G_{L,0}(x,t)=\frac12x_<\left(x_>-\frac1{x_>}\right),
\qquad G_{D,0}(x,t)=\frac1{2t}x_<\left(x_>-\frac1{x_>}\right).}
\]
These are also the $k\downarrow0$ limits of the previous expressions.
No nonzero multiple of $x$ vanishes at $1$, so this case is nonresonant.
For real negative $k$, replace it by $|k|$ throughout, since the printed
operator depends only on $k^2$.

\Needspace{5\baselineskip}
\par\bigskip
\noindent\textbf{10.1.6}\quad Given that
\[
L=(1-x^{2})\frac{d^{2}}{dx^{2}}-2x\frac{d}{dx}
\]
and that $G(\pm1,t)$ remains finite, show that no Green's function can be
constructed by the techniques of this section.

\emph{Note.}~The solutions to $L=0$ needed for the regions $x<t$ and $x>t$ are
linearly dependent.
\solution
Put the operator in divergence form:
\[
Ly=\frac d{dx}\bigl[(1-x^2)y'(x)\bigr].
\]
For a homogeneous solution, one integration gives
$(1-x^2)y'=C$, and a second gives
\[
y(x)=D+\frac C2\log\frac{1+x}{1-x}.
\]
The logarithm diverges at either endpoint. Thus a homogeneous solution
finite at $-1$ is a constant, and a homogeneous solution finite at $+1$
is also a constant.

For an interior source $-1<t<1$, the proposed Green's function would
therefore have to be
\[
G(x,t)=\begin{cases}A(t),&x<t,\\B(t),&x>t.\end{cases}
\]
It must be continuous at $t$, because a discontinuity in $G$ would produce
a delta derivative in $LG$, which is absent from the required source.
So $A(t)=B(t)$. Both ordinary derivatives are zero, and hence
\[
\left[(1-x^2)G_x(x,t)\right]_{t^-}^{t^+}=0.
\]
But integrating $LG=\delta(x-t)$ across $t$ requires this flux jump
to be $1$. This contradiction rules out an ordinary Green's function
under the specified finite-endpoint conditions.

The obstruction is the constant zero mode $L1=0$. For a regular source
$f$, endpoint finiteness of a solution to $Ly=f$ forces the endpoint
fluxes $(1-x^2)y'$ to vanish: a nonzero limiting flux would integrate
to the logarithmic divergence displayed above. Consequently
\[
\int_{-1}^1f(x)\,dx
=\left[(1-x^2)y'(x)\right]_{-1}^1=0.
\]
A delta source has integral $1$ and violates that condition. Even when
a source is compatible, an arbitrary additive constant remains.
A Green's function with the constant mode projected out would solve a
modified source equation, not $LG=\delta$ requested here.

\Needspace{5\baselineskip}
\par\bigskip
\noindent\textbf{10.1.7}\quad Find the Green's function for
\[
\frac{d^{2}\psi}{dt^{2}}+k\frac{d\psi}{dt}=f(t),
\]
subject to the initial conditions $\psi(0)=\psi'(0)=0$, and solve this ODE for
$t>0$ given $f(t)=\exp(-t)$.
\solution
Use $u$ for the source time and $t$ for the observation time.
For $t<u$, the zero initial data require the causal kernel to vanish.
For $t>u$, the homogeneous equation is $G_{tt}+kG_t=0$. When $k\ne0$
its solution is $A(u)+B(u)e^{-kt}$. Continuity at $t=u$ requires
$A+B e^{-ku}=0$, so it can be written $C(u)(1-e^{-k(t-u)})$.
Its first derivative at $u^+$ is $kC(u)$, while that at $u^-$ is zero.
The coefficient of $\psi''$ is one; therefore the jump must be one,
and $C(u)=1/k$. We obtain
\[
\boxed{G(t,u)=\Theta(t-u)\frac{1-e^{-k(t-u)}}k\quad(k\ne0).}
\]
For $k=0$, integrate $G_{tt}=0$ on each side and use the same matching
conditions, or take the limit:
\[
\boxed{G(t,u)=\Theta(t-u)(t-u)\quad(k=0).}
\]
For either expression $G$ is continuous with value zero at the source
and $[G_t]=1$, proving $(\partial_t^2+k\partial_t)G=\delta(t-u)$.

\erratum{The manual's kernel on pp.~119--120 is normalized for the
multiplied equation $(e^{kt}\psi')'=e^{kt}f$. It is $e^{-ku}G(t,u)$
and must be integrated against $e^{ku}f(u)$. The kernel above acts
directly on the original source $f(u)$. Also, the apparent singularities
at $k=0$ and $k=1$ in the final algebraic formula are removable and need
the cases given below.}

For $f(u)=e^{-u}$ and $k\ne0,1$, the convolution is
\begin{align*}
\psi(t)&=\frac1k\int_0^t\left(1-e^{-k(t-u)}\right)e^{-u}\,du\\
&=\frac1k\left[1-e^{-t}-e^{-kt}\int_0^t e^{(k-1)u}\,du\right]\\
&=\frac1k\left[1-e^{-t}-\frac{e^{-t}-e^{-kt}}{k-1}\right]\\
&=\boxed{\frac{(k-1)-ke^{-t}+e^{-kt}}{k(k-1)}}.
\end{align*}
Its derivatives are
\[
\psi'(t)=\frac{e^{-t}-e^{-kt}}{k-1},\qquad
\psi''(t)=\frac{-e^{-t}+ke^{-kt}}{k-1},
\]
so $\psi''+k\psi'=e^{-t}$ and $\psi(0)=\psi'(0)=0$.
For the exceptional values, evaluate the integrals directly:
\begin{align*}
k=1:\quad
\psi(t)&=\int_0^t(e^{-u}-e^{-t})\,du
=\boxed{1-(1+t)e^{-t}},\\
k=0:\quad
\psi(t)&=\int_0^t(t-u)e^{-u}\,du
=\boxed{t-1+e^{-t}}.
\end{align*}
For $k=1$, the derivatives are $te^{-t}$ and $(1-t)e^{-t}$;
for $k=0$, they are $1-e^{-t}$ and $e^{-t}$. These verify the ODE
and zero initial data in both cases. There is no boundary condition at
infinite time, so a negative real $k$ is also permitted in the causal
formula despite the possible growth of the response.

\Needspace{5\baselineskip}
\par\bigskip
\noindent\textbf{10.1.8}\quad \erratum{Take $k>0$ and time dependence
$e^{-i\omega\tau}$, with $\omega>0$.} Verify that the Green's function
\[
G(x,x')=-\frac{i}{2k}\exp\left(ik|x-x'|\right)
\]
yields an outgoing wave solution to the ODE
\[
\left(\frac{d^{2}}{dx^{2}}+k^{2}\right)\psi(x)=g(x).
\]
\emph{Note.}~Compare with Example 10.1.3.
\solution
\erratum{Take $k>0$ and the time convention $e^{-i\omega\tau}$ used
in Example 10.1.3. Its Eq.~(10.28), textbook p.~454, omits $k$ inside
the exponential; the expression in this exercise and Table 10.1 has
the correct factor $ik|x-x'|$.}
Let $s=x-x'$. Away from $s=0$, each branch is proportional to
$e^{iks}$ or $e^{-iks}$, and hence solves $G_{xx}+k^2G=0$.
Both branches have the common value $-i/(2k)$ at $s=0$. Their first
derivatives are
\[
G_x(x,x')=\begin{cases}
-\frac12e^{ik(x'-x)},&x<x',\\
\frac12e^{ik(x-x')},&x>x'.
\end{cases}
\]
Thus $G_x(x'^+,x')-G_x(x'^-,x')=1$. Continuity rules out a delta
derivative, and this jump gives
\[
(\partial_x^2+k^2)G(x,x')=\delta(x-x').
\]

For a continuous compactly supported source $g$, define
\[
F(x)=\int_{-\infty}^xe^{-ikt}g(t)\,dt,\qquad
H(x)=\int_x^\infty e^{ikt}g(t)\,dt.
\]
The convolution has the explicit form
\[
\psi(x)=-\frac{i}{2k}\left[e^{ikx}F(x)+e^{-ikx}H(x)\right].
\]
Because $F'=e^{-ikx}g$ and $H'=-e^{ikx}g$, their endpoint contributions
cancel in the first derivative. A second differentiation gives
\begin{align*}
\psi'(x)&=\frac12\left[e^{ikx}F(x)-e^{-ikx}H(x)\right],\\
\psi''(x)&=\frac{ik}{2}\left[e^{ikx}F(x)+e^{-ikx}H(x)\right]+g(x).
\end{align*}
Since $k^2\psi=-ik[e^{ikx}F+e^{-ikx}H]/2$, their sum is $g$.

If the source is supported in $[a,b]$, then for $x>b$ only the first
term remains: $\psi=C_+e^{ikx}$. For $x<a$ only the second remains:
$\psi=C_-e^{-ikx}$. With the time factor $e^{-i\omega\tau}$, constant
phase on the right obeys $kx-\omega\tau=\text{constant}$ and moves
to increasing $x$; on the left it obeys $-kx-\omega\tau=\text{constant}$
and moves to decreasing $x$. Both waves travel away from the source,
which verifies the outgoing condition.

\Needspace{5\baselineskip}
\par\bigskip
\noindent\textbf{10.1.9}\quad Construct the 1-D Green's function for the
modified Helmholtz equation, \erratum{with $k>0$ in the displayed answer},
\[
\left(\frac{d^{2}}{dx^{2}}-k^{2}\right)\psi(x)=f(x).
\]
The boundary conditions are that the Green's function must vanish for
$x\to\infty$ and $x\to-\infty$.
\par\nobreak\hfill\emph{ANS.}\enspace
$G(x_{1},x_{2})=-\dfrac{1}{2k}\exp\left(-k|x_{1}-x_{2}|\right)$.
\solution
\erratum{The printed formula assumes $k>0$. For a real nonzero
parameter of either sign, set $\kappa=|k|$; the operator contains
$k^2=\kappa^2$, while decay requires the positive root.}
For $x<t$, decay as $x\to-\infty$ selects $e^{\kappa x}$;
for $x>t$, decay as $x\to+\infty$ selects $e^{-\kappa x}$.
Continuity can be imposed by a single coefficient:
\[
G(x,t)=C\begin{cases}
e^{\kappa(x-t)},&x<t,\\
e^{-\kappa(x-t)},&x>t.
\end{cases}
\]
The jump of its derivative is $-\kappa C-\kappa C=-2\kappa C$.
The coefficient of $\psi''$ is one, so this must equal one. Hence
\[
\boxed{G(x,t)=-\frac1{2\kappa}e^{-\kappa|x-t|}.}
\]
It is continuous at $x=t$, solves $G_{xx}-\kappa^2G=0$ off the
diagonal, and decays at both infinities. Its derivative limits are
$-1/2$ and $+1/2$, so $(\partial_x^2-k^2)G=\delta(x-t)$.
Therefore
\[
\psi(x)=-\frac1{2\kappa}\left[
e^{-\kappa x}\int_{-\infty}^x e^{\kappa t}f(t)\,dt
+e^{\kappa x}\int_x^\infty e^{-\kappa t}f(t)\,dt\right]
\]
is the required solution for, for example, a continuous compactly
supported source. A homogeneous solution $Ae^{\kappa x}+Be^{-\kappa x}$
can decay at both ends only if $A=B=0$, proving uniqueness.

\erratum{For $k=0$, no Green's function for $d^2/dx^2$ can vanish at
both infinities. Each branch would be affine; decay at its corresponding
infinity would make that branch zero, leaving no derivative jump.}
The divergent constant in the $\kappa\downarrow0$ expansion
\[
G(x,t)=-\frac1{2\kappa}+\frac12|x-t|+O(\kappa)
\]
reflects that obstruction. Subtracting the constant would yield the
Laplace fundamental solution, but it would not satisfy decay at infinity.

\Needspace{5\baselineskip}
\par\bigskip
\noindent\textbf{10.1.10}\quad From the eigenfunction expansion of the Green's
function show that
\begin{enumerate}
\item[(a)]
  $\displaystyle\frac{2}{\pi^{2}}\sum_{n=1}^{\infty}
  \frac{\sin n\pi x\,\sin n\pi t}{n^{2}}
  =\begin{cases}
  x(1-t), & 0\le x<t,\\
  t(1-x), & t<x\le 1.
  \end{cases}$
\item[(b)]
  $\displaystyle\frac{2}{\pi^{2}}\sum_{n=0}^{\infty}
  \frac{\sin(n+\frac{1}{2})\pi x\,\sin(n+\frac{1}{2})\pi t}
  {(n+\frac{1}{2})^{2}}
  =\begin{cases}
  x, & 0\le x<t,\\
  t, & t<x\le 1.
  \end{cases}$
\end{enumerate}
\solution
In both parts use $L=-d^2/dx^2$ on $[0,1]$. For a normalized
eigenfunction $\phi_n$, $L\phi_n=\lambda_n\phi_n$, the inverse has
coefficients $1/\lambda_n$:
\[
G(x,t)=\sum_n\frac{\phi_n(x)\phi_n(t)}{\lambda_n}.
\]
There is no complex conjugation to display here because the chosen
eigenfunctions are real.

(a) With $\phi(0)=\phi(1)=0$, the equation $-\phi''=\lambda\phi$
has positive eigenvalues. Indeed, integration by parts gives
\[
\lambda\int_0^1|\phi|^2dx
=\int_0^1|\phi'|^2dx,
\]
and equality at $\lambda=0$ would make $\phi$ a constant vanishing
at the endpoints. Write $\lambda=q^2$ with $q>0$. The condition at zero
removes the cosine term, and $\sin q=0$ gives $q=n\pi$, $n\ge1$.
Since
\[
\int_0^1\sin^2(n\pi x)\,dx
=\left[\frac x2-\frac{\sin(2n\pi x)}{4n\pi}\right]_0^1=\frac12,
\]
we have $\phi_n(x)=\sqrt2\sin(n\pi x)$ and $\lambda_n=n^2\pi^2$.
The eigenfunction formula is therefore
\[
G_D(x,t)=\frac2{\pi^2}\sum_{n=1}^\infty
\frac{\sin(n\pi x)\sin(n\pi t)}{n^2}.
\]
Construct the same kernel directly. The left and right homogeneous
solutions are $x$ and $1-x$, and continuity suggests
$C x_<(1-x_>)$. Its derivative jump is
$-Ct-C(1-t)=-C$. Requiring $-G_{xx}=\delta$ gives $C=1$.
Thus
\[
\boxed{\frac2{\pi^2}\sum_{n=1}^\infty
\frac{\sin(n\pi x)\sin(n\pi t)}{n^2}
=\begin{cases}x(1-t),&x\le t,\\t(1-x),&t\le x.\end{cases}}
\]

(b) Now impose $\phi(0)=0$, $\phi'(1)=0$. The same energy identity
shows that the eigenvalues are positive: its boundary term vanishes
also for these mixed conditions. Again $\phi=A\sin(qx)$, while the
second endpoint condition gives $q\cos q=0$. Thus
\[
q_n=(n+\tfrac12)\pi,\qquad
\lambda_n=(n+\tfrac12)^2\pi^2,\qquad n=0,1,2,\ldots.
\]
The normalization is unchanged because $\sin(2q_n)=0$:
$\int_0^1\sin^2(q_nx)dx=1/2$. Hence
\[
G_M(x,t)=\frac2{\pi^2}\sum_{n=0}^\infty
\frac{\sin((n+\tfrac12)\pi x)\sin((n+\tfrac12)\pi t)}{(n+\tfrac12)^2}.
\]
Exercise 10.1.1 constructed the unique kernel for precisely these
conditions, $G_M(x,t)=\min(x,t)$. It follows that
\[
\boxed{\frac2{\pi^2}\sum_{n=0}^\infty
\frac{\sin((n+\tfrac12)\pi x)\sin((n+\tfrac12)\pi t)}{(n+\tfrac12)^2}
=\begin{cases}x,&x\le t,\\t,&t\le x.\end{cases}}
\]

For completeness, the identification does not require differentiating
either series twice term by term. For each fixed interior $t$, integrate
by parts on $(0,t)$ and $(t,1)$, using $-G_{xx}=\delta(x-t)$ and
the matching endpoint conditions.  Because each closed kernel is affine
on either side of $t$, integration by parts twice gives
\begin{align*}
\int_0^1G(-\phi_n'')\,dx
&=\bigl[-G\phi_n'+G_x\phi_n\bigr]_{0}^{t^-}
 +\bigl[-G\phi_n'+G_x\phi_n\bigr]_{t^+}^{1}\\
&=\bigl[G_x(t^-,t)-G_x(t^+,t)\bigr]\phi_n(t)=\phi_n(t).
\end{align*}
Here the outer-endpoint terms vanish; the terms with $\phi_n'(t)$
cancel because $G$ is continuous, and the derivative jump is $-1$.
Since $-\phi_n''=\lambda_n\phi_n$, this gives
\[
\lambda_n\int_0^1G(x,t)\phi_n(x)\,dx
=\int_0^1G(x,t)(-\phi_n''(x))\,dx=\phi_n(t).
\]
Thus the closed kernels have exactly the stated Fourier coefficients.
Completeness gives equality in $L^2$. Both series converge absolutely
and uniformly by comparison with $\sum n^{-2}$ or
$\sum(n+1/2)^{-2}$, so their sums are continuous. Equality therefore
holds everywhere, including $x=t$ and the endpoints by continuity.

\Needspace{5\baselineskip}
\par\bigskip
\noindent\textbf{10.1.11}\quad Derive an integral equation corresponding to
\[
y''(x)-y(x)=0,\quad y(1)=1,\quad y(-1)=1,
\]
\begin{enumerate}
\item[(a)] by integrating twice.
\item[(b)] by forming the Green's function.
\end{enumerate}
\par\nobreak\hfill\emph{ANS.}\enspace
$y(x)=1-\displaystyle\int_{-1}^{1}K(x,t)\,y(t)\,dt$,\par
\nobreak\hfill
$K(x,t)=\begin{cases}
\frac{1}{2}(1-x)(t+1), & x>t,\\
\frac{1}{2}(1-t)(x+1), & x<t.
\end{cases}$
\solution
(a) Rewrite the ODE as $y''=y$, and set $c=y'(-1)$.
Integrating first to $s$ and then to $x$ gives
\[
y'(s)=c+\int_{-1}^s y(t)\,dt,
\qquad
y(x)=1+c(x+1)+\int_{-1}^x\int_{-1}^s y(t)\,dt\,ds.
\]
Interchanging the order on the triangle $-1\le t\le s\le x$ turns
the double integral into $\int_{-1}^x(x-t)y(t)dt$. At $x=1$,
\[
1=1+2c+\int_{-1}^1(1-t)y(t)\,dt,
\qquad c=-\frac12\int_{-1}^1(1-t)y(t)\,dt.
\]
Substitution yields
\[
y(x)=1+\int_{-1}^x(x-t)y(t)\,dt
-\frac{x+1}{2}\int_{-1}^1(1-t)y(t)\,dt.
\]
Split the last integral at $t=x$. On $t<x$, its combined coefficient is
\[
(x-t)-\frac{x+1}{2}(1-t)
=\frac{2x-2t-x-1+xt+t}{2}
=-\frac{(1-x)(t+1)}2.
\]
On $t>x$, the coefficient is $-(x+1)(1-t)/2$. We obtain
\[
\boxed{y(x)=1-\int_{-1}^1K(x,t)y(t)\,dt,\quad
K(x,t)=\begin{cases}
\frac12(1-x)(t+1),&t\le x,\\
\frac12(1-t)(x+1),&x\le t.
\end{cases}}
\]

(b) Homogenize the boundary conditions by setting $u=y-1$.
Then $u(-1)=u(1)=0$, and the original equation gives
$-u''=-y$. Construct the Green's function of $-d^2/dx^2$ on $[-1,1]$:
the left and right homogeneous solutions are $x+1$ and $1-x$.
The continuous kernel $C(x_<+1)(1-x_>)$ has derivative jump
\[
-C(t+1)-C(1-t)=-2C.
\]
Setting that jump to $-1$ gives $C=1/2$, which is exactly $K$ above.
Thus
\[
u(x)=\int_{-1}^1K(x,t)(-y(t))\,dt,
\]
and adding $1$ recovers the same integral equation.

Conversely, $K(\pm1,t)=0$ makes any solution of that integral equation
satisfy $y(\pm1)=1$. Applying $-d^2/dx^2$ gives $-y''=-y$, or $y''-y=0$.
As an explicit check, write $y=A\cosh x+B\sinh x$. The two endpoint
values give $A\cosh1+B\sinh1=1$ and $A\cosh1-B\sinh1=1$,
so
\[
\boxed{y(x)=\frac{\cosh x}{\cosh1}.}
\]
Direct integration of the two branches of $K$ gives
$\int_{-1}^1K(x,t)\cosh t\,dt=\cosh1-\cosh x$, verifying the
integral equation with this solution.

\Needspace{5\baselineskip}
\par\bigskip
\noindent\textbf{10.1.12}\quad The general second-order linear ODE with
constant coefficients is
\[
y''(x)+a_{1}y'(x)+a_{2}y(x)=0.
\]
Given the boundary conditions $y(0)=y(1)=0$, integrate twice and develop the
integral equation
\[
y(x)=\int_{0}^{1}K(x,t)\,y(t)\,dt,
\]
with
\[
K(x,t)=\begin{cases}
a_{2}t(1-x)+a_{1}(x-1), & t<x,\\
a_{2}x(1-t)+a_{1}x, & x<t.
\end{cases}
\]
Note that $K(x,t)$ is symmetric and continuous if $a_{1}=0$. How is this
related to self-adjointness of the ODE?
\erratum{For this self-adjointness discussion take $a_1,a_2$ real,
and distinguish the ordinary unit-weight inner product from a weighted one.}
\solution
Let $c=y'(0)$ and use $y(0)=0$. One integration of
$y''=-a_1y'-a_2y$ gives
\[
y'(x)=c-a_1y(x)-a_2\int_0^xy(t)\,dt.
\]
Integrate once more. Reversing the order in the double integral gives
\[
y(x)=cx-a_1\int_0^xy(t)\,dt-a_2\int_0^x(x-t)y(t)\,dt.
\]
The second boundary condition supplies
\[
0=c-a_1\int_0^1y(t)\,dt-a_2\int_0^1(1-t)y(t)\,dt.
\]
Eliminate $c$ from the preceding expression:
\begin{align*}
y(x)={}&a_1x\int_0^1y(t)\,dt-a_1\int_0^xy(t)\,dt\\
&+a_2x\int_0^1(1-t)y(t)\,dt-a_2\int_0^x(x-t)y(t)\,dt.
\end{align*}
For $t<x$, the coefficient of $y(t)$ is
\[
a_1(x-1)+a_2\bigl[x(1-t)-(x-t)\bigr]
=a_1(x-1)+a_2t(1-x).
\]
For $t>x$, it is $a_1x+a_2x(1-t)$. Hence
\[
\boxed{y(x)=\int_0^1K(x,t)y(t)\,dt,\qquad
K(x,t)=\begin{cases}
a_2t(1-x)+a_1(x-1),&t<x,\\
a_2x(1-t)+a_1x,&t>x.
\end{cases}}
\]
The value of $K$ at the single point $t=x$ does not affect the integral.

The reverse implication can also be checked without assuming the ODE.
Let the integral on the right be $F(x)$ and split it at $t=x$. The
first branch at the diagonal minus the second branch there is $-a_1$,
so the Leibniz rule gives
\[
F'(x)=-a_1y(x)+\int_0^x(-a_2t+a_1)y(t)\,dt
+\int_x^1\bigl[a_2(1-t)+a_1\bigr]y(t)\,dt.
\]
A second differentiation gives
\[
F''=-a_1y'+\bigl[(-a_2x+a_1)-(a_2(1-x)+a_1)\bigr]y
=-a_1y'-a_2y.
\]
Thus $y=F$ implies the original ODE. Also $K(0,t)=K(1,t)=0$,
so the integral equation incorporates both boundary conditions.

\erratum{The self-adjointness discussion assumes real $a_1,a_2$ and
the unit-weight inner product. A nonzero real $a_1$ does not make the
ODE intrinsically incapable of a self-adjoint formulation; it requires
an integrating factor, or equivalently a weighted inner product.}
For real coefficients and zero Dirichlet data, integration by parts gives
\[
\left(\frac{d^2}{dx^2}+a_1\frac d{dx}+a_2\right)^{\!\dagger}
=\frac{d^2}{dx^2}-a_1\frac d{dx}+a_2.
\]
So the displayed operator is formally self-adjoint in unit weight
precisely when $a_1=0$. In this case
\[
K(x,t)=a_2x_<(1-x_>),
\]
which is continuous and symmetric. When $a_1\ne0$, the kernel jumps by
$a_1$ as $t$ passes from below to above $x$ and is not symmetric.
Multiplication of the ODE by $p(x)=e^{a_1x}$ instead gives
\[
(py')'+a_2py=0.
\]
This differential expression is self-adjoint with the Dirichlet domain;
the original expression is symmetric in the weighted scalar product
$\int_0^1p(x)\phi(x)^*\psi(x)dx$.

The homogeneous integral equation need not have a nonzero solution for
arbitrary coefficients. To identify its exceptional parameters, put
$y=e^{-a_1x/2}v$. Then
\[
y'=e^{-a_1x/2}(v'-a_1v/2),\qquad
y''=e^{-a_1x/2}(v''-a_1v'+a_1^2v/4),
\]
and substitution reduces the ODE to
\[
v''+(a_2-a_1^2/4)v=0,\qquad v(0)=v(1)=0.
\]
For real coefficients, nonzero solutions exist exactly when
$a_2-a_1^2/4=n^2\pi^2$, $n=1,2,\ldots$; they are
$y=C e^{-a_1x/2}\sin(n\pi x)$. Otherwise the only solution of either
formulation is $y=0$.

\Needspace{5\baselineskip}
\par\bigskip
\noindent\textbf{10.1.13}\quad Transform the ODE
\[
\frac{d^{2}y(r)}{dr^{2}}-k^{2}y(r)+V_{0}\frac{e^{-r}}{r}y(r)=0
\]
and the boundary conditions $y(0)=y(\infty)=0$ into an integral equation of
the form
\[
y(r)=-V_{0}\int_{0}^{\infty}G(r,t)\,\frac{e^{-t}}{t}\,y(t)\,dt.
\]
The quantities $V_{0}$ and $k^{2}$ are constants.
\erratum{Take $k>0$ (or, more generally, $\Re k>0$ for the displayed
decaying kernel), and the regular branch $y(r)=O(r)$ at the origin.}
The ODE is derived from the
Schr\"odinger wave equation with a mesonic potential:
\[
G(r,t)=\begin{cases}
-\dfrac{1}{k}e^{-kt}\sinh kr, & 0\le r<t,\\[6pt]
-\dfrac{1}{k}e^{-kr}\sinh kt, & t<r<\infty.
\end{cases}
\]
\solution
\erratum{The decaying kernel in the question assumes $k$ is the
positive square root of a positive $k^2$ (more generally $\Re k>0$).
The singular endpoint is interpreted using the regular solution
$y(r)=O(r)$ as $r\downarrow0$. The manual's final sentence on p.~120,
claiming that the question's $G$ is minus the Green's function, is
incorrect for the chapter's convention $LG=+\delta$.}
Isolate the free operator and source:
\[
L_0=\frac{d^2}{dr^2}-k^2,\qquad
L_0y(r)=-V_0\frac{e^{-r}}r y(r).
\]
For $L_0u=0$, choose $u(r)=\sinh(kr)$, which vanishes at zero,
and $v(r)=e^{-kr}$, which decays at infinity. Their Wronskian is
\begin{align*}
uv'-u'v
&=-k\sinh(kr)e^{-kr}-k\cosh(kr)e^{-kr}\\
&=-ke^{-kr}\bigl[\sinh(kr)+\cosh(kr)\bigr]=-k.
\end{align*}
The continuous kernel $C u(r_<)v(r_>)$ has derivative jump $-kC$.
For $L_0G=\delta(r-t)$ the jump must be $1$, so $C=-1/k$:
\[
\boxed{G(r,t)=-\frac1k\sinh(kr_<)e^{-kr_>},
\qquad r_<=\min(r,t),\quad r_>=\max(r,t).}
\]
It satisfies $G(0,t)=0$ and $G(\infty,t)=0$. The matching calculation
shows directly that it is the Green's function with a positive unit
delta source. Equivalently, the method-of-images form is
\[
G(r,t)=-\frac1{2k}\left[e^{-k|r-t|}-e^{-k(r+t)}\right].
\]
The second term cancels the first at $r=0$ and is homogeneous as a
function of positive $r$.

Apply this inverse to the source $-V_0e^{-t}y(t)/t$:
\[
\boxed{y(r)=-V_0\int_0^\infty G(r,t)\frac{e^{-t}}t\,y(t)\,dt.}
\]
Thus the negative sign outside the integral and the negative sign in
$G$ are both required. In split form, the same equation reads
\[
y(r)=\frac{V_0}{k}\left[
e^{-kr}\int_0^r\sinh(kt)\frac{e^{-t}}t y(t)\,dt
+\sinh(kr)\int_r^\infty e^{-kt}\frac{e^{-t}}t y(t)\,dt\right].
\]
On the regular branch $y(t)=O(t)$, the source factor $y(t)/t$ is bounded
near zero. For example, substitution of $y=At+Bt^2+\cdots$ into the
ODE gives $2B+V_0A=0$, so
$y(t)=At-(V_0A/2)t^2+\cdots$. At infinity the exponential factors
ensure convergence for a decaying $y$.

Conversely, apply $L_0$ to the integral equation. The delta identity
gives $L_0y=-V_0e^{-r}y/r$, and the kernel enforces the endpoint
conditions. There is no omitted homogeneous term: $Ae^{kr}+Be^{-kr}$
can decay at infinity only if $A=0$, and its value at zero then forces
$B=0$. The integral equation is therefore equivalent to the original
boundary problem within this regular, decaying class; it is a homogeneous
eigenvalue equation for the allowed parameters, not an assertion that
all $V_0,k$ admit a nonzero solution.

The limit $k\downarrow0$ gives $G(r,t)\to-\min(r,t)$, which approaches
$-t$, not zero, as $r\to\infty$. Accordingly there is no inverse of
$d^2/dr^2$ satisfying these same two zero boundary values for an arbitrary
source: its right homogeneous branch would have to be identically zero.
This explains why the decay assumption on the free parameter is needed.

\Needspace{9\baselineskip}
\subsection*{10.2 Problems in Two and Three Dimensions}

\Needspace{5\baselineskip}
\par\bigskip
\noindent\textbf{10.2.1}\quad Show that the fundamental Green's function for
the 1-D Laplace equation, $|x_{1}-x_{2}|/2$, is consistent with the form found
in Example 10.1.1.
\solution
Write $x=x_1$, $t=x_2$. The fundamental solution
$G_0(x,t)=|x-t|/2$ has derivative $-1/2$ for $x<t$ and $+1/2$ for
$x>t$. Its derivative jump is $1$, so
\[
\partial_x^2G_0(x,t)=\delta(x-t).
\]
To impose zero Dirichlet values at $x=0,1$, add a homogeneous solution
in $x$, $h(x,t)=a(t)x+b(t)$. The endpoint equations are
\[
0=G_0(0,t)+b(t)=t/2+b(t),
\]
\[
0=G_0(1,t)+a(t)+b(t)=(1-t)/2+a(t)-t/2.
\]
Hence $b=-t/2$, $a=t-1/2$, and
\[
G_D(x,t)=\frac12|x-t|+xt-\frac{x+t}{2}.
\]
This added affine function has zero second derivative and cannot change
the source normalization. Evaluating the two branches gives
\begin{align*}
x<t:\quad G_D&=\frac{t-x}{2}+xt-\frac{x+t}{2}
=-x(1-t),\\
x>t:\quad G_D&=\frac{x-t}{2}+xt-\frac{x+t}{2}
=-t(1-x).
\end{align*}
Thus
\[
\boxed{G_D(x,t)=-\begin{cases}x(1-t),&x\le t,\\t(1-x),&t\le x.\end{cases}}
\]
\erratum{The operator in Example 10.1.1 is $-d^2/dx^2$, whereas the
Laplace operator here is $+d^2/dx^2$. Their Green's functions therefore
differ by a minus sign. Consistency requires retaining that operator
sign, as also shown in Example 10.2.1, textbook p.~464.}
Indeed $\partial_x^2G_D=\delta$, while
$(-\partial_x^2)(-G_D)=\delta$; the latter positive kernel is exactly
the one in Example 10.1.1.

\Needspace{5\baselineskip}
\par\bigskip
\noindent\textbf{10.2.2}\quad Show that if
\[
L\psi(\mathbf{r})\equiv
\nabla\cdot\left[p(\mathbf{r})\nabla\psi(\mathbf{r})\right]
+q(\mathbf{r})\psi(\mathbf{r}),
\]
then $L$ is Hermitian for $p(\mathbf{r})$ and $q(\mathbf{r})$ real, assuming
\erratum{common homogeneous Dirichlet boundary conditions, sufficient
regularity for Green's identity, and a vanishing boundary form at infinity
if the region is unbounded}, and that the scalar product is an integral
over that region with unit weight.
\solution
\erratum{For an operator domain, the Dirichlet conditions used here
are homogeneous: both test functions vanish at the boundary.
Nonzero prescribed boundary data must first be subtracted using a fixed
extension. On an unbounded region the corresponding boundary form at
infinity must also vanish.}
Let $\Omega$ have a smooth boundary, with sufficiently regular
$p,q,\phi,\psi$ for the integrations below, and take $p,q$ real.
The unit-weight scalar product is
\[
\langle\phi,\psi\rangle=\int_\Omega\phi^*\psi\,dV.
\]
The product rule gives
\[
\phi^*\nabla\cdot(p\nabla\psi)
=\nabla\cdot(p\phi^*\nabla\psi)-p\nabla\phi^*\cdot\nabla\psi.
\]
Integrate and use the divergence theorem:
\begin{align*}
\langle\phi,L\psi\rangle
&=\int_{\partial\Omega}p\phi^*\partial_n\psi\,dS
-\int_\Omega p\nabla\phi^*\cdot\nabla\psi\,dV
+\int_\Omega q\phi^*\psi\,dV\\
&=-\int_\Omega p\nabla\phi^*\cdot\nabla\psi\,dV
+\int_\Omega q\phi^*\psi\,dV,
\end{align*}
because $\phi=0$ on $\partial\Omega$.
Reality of the coefficients gives
$(L\phi)^*=\nabla\cdot(p\nabla\phi^*)+q\phi^*$. Repeating the
integration by parts with $\psi$ outside the divergence yields
\begin{align*}
\langle L\phi,\psi\rangle
&=\int_{\partial\Omega}p\psi\partial_n\phi^*\,dS
-\int_\Omega p\nabla\psi\cdot\nabla\phi^*\,dV
+\int_\Omega q\phi^*\psi\,dV\\
&=-\int_\Omega p\nabla\psi\cdot\nabla\phi^*\,dV
+\int_\Omega q\phi^*\psi\,dV.
\end{align*}
The remaining volume integrals agree term by term, so
\[
\boxed{\langle\phi,L\psi\rangle=\langle L\phi,\psi\rangle.}
\]
Equivalently, their difference is the boundary form
\[
\int_{\partial\Omega}p\left(\phi^*\partial_n\psi
-\psi\partial_n\phi^*\right)dS=0.
\]
This proves the Hermitian (symmetric) property on the stated common
domain. \erratum{If ``self-adjoint'' is used in its stronger operator-theoretic
sense, a domain realization must also be specified; the usual smooth,
bounded, uniformly elliptic Dirichlet realization is one such setting.}

\Needspace{5\baselineskip}
\par\bigskip
\noindent\textbf{10.2.3}\quad Show that the terms $+k^{2}$ in the Helmholtz
operator and $-k^{2}$ in the modified Helmholtz operator do not affect the
behavior of $G(\mathbf{r}_{1},\mathbf{r}_{2})$ in the immediate vicinity of
the singular point $\mathbf{r}_{1}=\mathbf{r}_{2}$. Specifically, show that
\[
\erratum{\lim_{\varepsilon\downarrow0}
\int_{|\mathbf{r}_{2}-\mathbf{r}_{1}|<\varepsilon}
k^{2}G(\mathbf{r}_{1},\mathbf{r}_{2})\,d^{3}r_{2}=0.}
\]
\solution
\erratum{The printed limit, including its value $-1$, is incorrect
for a shrinking neighborhood of the source (textbook p.~466). The
integration region must be specified as a ball of radius $\varepsilon$,
and the local limit is $0$. The manual, p.~121, computes a different,
whole-space Fourier integral, formally giving $+1/k^2$ for its
Helmholtz denominator $k^2-p^2$.  It does not establish the printed
local claim.}
Fix $\mathbf r_1$ and put $s=|\mathbf r_2-\mathbf r_1|$.
The three-dimensional outgoing and modified kernels are respectively
\[
G_H(s)=-\frac{e^{iks}}{4\pi s},\qquad
G_M(s)=-\frac{e^{-ks}}{4\pi s},\qquad k>0.
\]
Their Taylor expansions at $s=0$ start with
\[
G_H(s)=-\frac1{4\pi s}-\frac{ik}{4\pi}+O(s),\qquad
G_M(s)=-\frac1{4\pi s}+\frac{k}{4\pi}+O(s).
\]
Both therefore have the same leading singularity as the Laplace kernel.
Since $d^3r_2=4\pi s^2ds$ for a radial integral,
\begin{align*}
\int_{s<\varepsilon}k^2G(s)\,d^3r_2
&=4\pi k^2\int_0^\varepsilon
\left[-\frac1{4\pi s}+O(1)\right]s^2\,ds\\
&=-\frac{k^2\varepsilon^2}{2}+O(\varepsilon^3).
\end{align*}
Consequently the corrected assertion is
\[
\boxed{\lim_{\varepsilon\downarrow0}
\int_{|\mathbf r_2-\mathbf r_1|<\varepsilon}
k^2G(\mathbf r_1,\mathbf r_2)\,d^3r_2=0.}
\]
The same conclusion holds for $-k^2G$.

The exact radial integrals independently show the limit. For Helmholtz,
\begin{align*}
I_H(\varepsilon)&=-k^2\int_0^\varepsilon s e^{iks}\,ds\\
&=-k^2\left[e^{iks}\left(\frac{s}{ik}+\frac1{k^2}\right)\right]_0^\varepsilon\\
&=e^{ik\varepsilon}(ik\varepsilon-1)+1\longrightarrow0.
\end{align*}
For modified Helmholtz,
\begin{align*}
I_M(\varepsilon)&=-k^2\int_0^\varepsilon s e^{-ks}\,ds\\
&=\left[(1+ks)e^{-ks}\right]_0^\varepsilon
=(1+k\varepsilon)e^{-k\varepsilon}-1\longrightarrow0.
\end{align*}
Integrating the defining equation over the ball now gives
\[
\int_{s=\varepsilon}\nabla G\cdot d\mathbf S
\ \mathord{\pm}\ \int_{s<\varepsilon}k^2G\,dV=1.
\]
The second integral tends to zero, while the flux tends to one.
Therefore it is the Laplacian term that supplies the delta singularity;
the zeroth-order $\pm k^2$ term cannot change its coefficient.
For comparison, taking the entire space for the modified kernel means
$\varepsilon\to\infty$, and gives $I_M\to-1$. That global value must
not be substituted for the shrinking-ball limit.

\Needspace{5\baselineskip}
\par\bigskip
\noindent\textbf{10.2.4}\quad Show that
\[
-\frac{\exp(ik|\mathbf{r}_{1}-\mathbf{r}_{2}|)}
{4\pi|\mathbf{r}_{1}-\mathbf{r}_{2}|}
\]
satisfies the appropriate criteria and therefore is a Green's function for the
Helmholtz equation. \erratum{Use $k>0$ and the outgoing convention with
time factor $e^{-i\omega\tau}$, $\omega>0$.}
\solution
Let $s=|\mathbf r_1-\mathbf r_2|$ and take $k>0$, with time
dependence $e^{-i\omega\tau}$. For a radial function $G(s)$ at $s>0$,
\[
\nabla^2G=G''+\frac2sG'.
\]
Write $G=F(s)/s$. Then
\[
G'=\frac{F'}s-\frac F{s^2},\qquad
G''=\frac{F''}s-\frac{2F'}{s^2}+\frac{2F}{s^3},
\]
so the first-derivative terms cancel and
\[
(\nabla^2+k^2)G=\frac{F''+k^2F}{s}.
\]
The homogeneous radial solutions are consequently
$G=(Ae^{iks}+Be^{-iks})/s$. The $e^{iks-i\omega\tau}$ phase moves
outward, whereas $e^{-iks-i\omega\tau}$ moves inward. Thus the outgoing
condition selects $B=0$.

For $G=Ae^{iks}/s$, the flux through a sphere of radius $\varepsilon$
centered on the source is
\[
4\pi\varepsilon^2G'(\varepsilon)
=4\pi A e^{ik\varepsilon}(ik\varepsilon-1)
\longrightarrow-4\pi A.
\]
Exercise 10.2.3 shows that the integral of $k^2G$ over that small ball
vanishes. A positive unit delta source therefore requires $-4\pi A=1$,
giving
\[
\boxed{G(\mathbf r_1,\mathbf r_2)
=-\frac{e^{ik|\mathbf r_1-\mathbf r_2|}}
{4\pi|\mathbf r_1-\mathbf r_2|}.}
\]
It solves the homogeneous Helmholtz equation away from the source, is
smooth there, and has a locally integrable $1/s$ singularity.
Its outgoing Sommerfeld condition is explicit:
\[
s(\partial_s-ik)G=\frac{e^{iks}}{4\pi s}\longrightarrow0.
\]

To verify the delta identity against an arbitrary smooth compactly
supported test function $\eta$, apply Green's second identity outside a
ball of radius $\varepsilon$ around the source. The outer boundary
contributes zero, while the outward normal of this punctured region
at the inner sphere points in the direction $-\widehat{\mathbf s}$.
Since $(\nabla^2+k^2)G=0$ there, the identity reads
\[
\int_{s>\varepsilon}G(\nabla^2+k^2)\eta\,dV
=\int_{s=\varepsilon}
\left[-G\,\partial_s\eta+\eta\,\partial_sG\right]dS.
\]
The first surface term is $O(\varepsilon)$, because $G=O(1/\varepsilon)$
and the sphere has area $O(\varepsilon^2)$. In the second term,
$\eta=\eta(\mathbf r_2)+O(\varepsilon)$ and the flux tends to one,
so its limit is $\eta(\mathbf r_2)$. This proves, with its sign and
normalization,
\[
(\nabla_1^2+k^2)G(\mathbf r_1,\mathbf r_2)
=\delta^{(3)}(\mathbf r_1-\mathbf r_2).
\]
Convolution with a compactly supported source therefore produces the
outgoing solution of the inhomogeneous Helmholtz equation.

\Needspace{5\baselineskip}
\par\bigskip
\noindent\textbf{10.2.5}\quad Find the Green's function for the 3-D Helmholtz
equation, Exercise 10.2.4, when the wave is a standing wave.
\erratum{Find the customary equal average of outgoing and incoming
kernels, and state the remaining freedom to add a regular standing
homogeneous solution.}
\solution
With $s=|\mathbf r_1-\mathbf r_2|$, the outgoing and incoming
Green's functions are
\[
G_{\rm out}=-\frac{e^{iks}}{4\pi s},\qquad
G_{\rm in}=-\frac{e^{-iks}}{4\pi s}.
\]
For real $k$, they are complex conjugates. Each solves
$(\nabla^2+k^2)G=\delta^{(3)}$; its small-sphere flux tends to $1$.
The equal average therefore still has a unit, rather than doubled,
delta source:
\begin{align*}
G_{\rm st}&=\frac12(G_{\rm out}+G_{\rm in})\\
&=-\frac{e^{iks}+e^{-iks}}{8\pi s}
=\boxed{-\frac{\cos(ks)}{4\pi s}}.
\end{align*}
Multiplying by $e^{-i\omega\tau}$ and taking the real part gives
$-\cos(ks)\cos(\omega\tau)/(4\pi s)$, a standing radial wave.
Equivalently, the outgoing and incoming amplitudes are equal, so their
opposite radial fluxes cancel. Directly, $\operatorname{Im}(G_{\rm st}^*
\partial_sG_{\rm st})=0$ because $G_{\rm st}$ is real.

Its derivative and source flux are
\[
G_{\rm st}'=\frac{k\sin(ks)}{4\pi s}
+\frac{\cos(ks)}{4\pi s^2},\qquad
4\pi s^2G_{\rm st}'=ks\sin(ks)+\cos(ks)\longrightarrow1.
\]
Also $F(s)=sG_{\rm st}=-\cos(ks)/(4\pi)$ satisfies $F''+k^2F=0$.
These checks establish the same differential operator and source as in
Exercise 10.2.4.

\erratum{A standing-wave condition alone does not uniquely specify the
regular homogeneous part. The answer above uses the customary equal
average of the outgoing and incoming Green's functions. For any real
constant $C$, adding $C\sin(ks)/s$ gives another real radial standing-wave
kernel with the same delta source; an additional phase convention or
boundary condition is needed to exclude it.}
Indeed $\sin(ks)/s=k-k^3s^2/6+\cdots$ is smooth at the origin,
its small-sphere flux tends to zero, and it satisfies the homogeneous
Helmholtz equation everywhere. Thus it does not alter the singular
normalization.

\Needspace{5\baselineskip}
\par\bigskip
\noindent\textbf{10.2.6}\quad Verify that the formula given for the 3-D
Green's function of the modified Helmholtz equation in Table 10.1 is correct
when the boundary conditions of the problem are that $G$ vanish at infinity.
\erratum{Use the positive decay root $k>0$ in the table and also identify
its $k\downarrow0$ limit.}
\solution
The three-dimensional entry in Table 10.1, textbook p.~463, is
\[
\boxed{G(\mathbf r_1,\mathbf r_2)=-\frac{e^{-ks}}{4\pi s},
\qquad s=|\mathbf r_1-\mathbf r_2|.}
\]
\erratum{Here $k>0$ is the decay constant. For a real nonzero parameter
of either sign, replace $k$ by $|k|$. The limit $k=0$ gives the decaying
three-dimensional Laplace kernel $-1/(4\pi s)$.}
For $s>0$, put $F=sG=-e^{-ks}/(4\pi)$. The radial calculation of
Exercise 10.2.4 gives
\[
(\nabla^2-k^2)G=\frac{F''-k^2F}{s}=0.
\]
The kernel decays at infinity, and its leading singularity is
$-1/(4\pi s)$.

Check the entire source normalization on a sphere of arbitrary radius
$R>0$, not just in a limiting expansion. The derivative is
\[
G'(s)=\frac{e^{-ks}}{4\pi}\left(\frac{k}s+\frac1{s^2}\right),
\]
so its flux is
\[
\Phi(R)=4\pi R^2G'(R)=(1+kR)e^{-kR}.
\]
The integral of $G$ over the ball is
\begin{align*}
I(R)&=4\pi\int_0^R G(s)s^2\,ds=-\int_0^R s e^{-ks}\,ds\\
&=\left[e^{-ks}\left(\frac s{k}+\frac1{k^2}\right)\right]_0^R
=\frac{(1+kR)e^{-kR}-1}{k^2}.
\end{align*}
Therefore
\[
\int_{s<R}(\nabla^2-k^2)G\,dV
=\Phi(R)-k^2I(R)
=(1+kR)e^{-kR}-\bigl[(1+kR)e^{-kR}-1\bigr]=1.
\]
In particular $\Phi(R)\to1$ and $k^2I(R)\to0$ as $R\downarrow0$.
The test-function argument from Exercise 10.2.4 applies without change
to the singularity, proving
\[
(\nabla_1^2-k^2)G=\delta^{(3)}(\mathbf r_1-\mathbf r_2).
\]

For uniqueness in the class of kernels with that singularity and decay,
the difference $h$ of two such solutions is regular everywhere,
vanishes at infinity, and obeys $\nabla^2h=k^2h$.
Apply this to its real part: a positive maximum would be attained at
a finite point and would have $\nabla^2h\le0$, contradicting $k^2h>0$;
a negative minimum gives the opposite contradiction. Hence the real
part is zero, and the same argument applies to the imaginary part.
This establishes uniqueness for $k>0$ under these conditions.

\Needspace{5\baselineskip}
\par\bigskip
\noindent\textbf{10.2.7}\quad An electrostatic potential (mks units) is
\[
\varphi(\mathbf{r})=Z\,\frac{e^{-ar}}{4\pi\epsilon_{0}r},
\qquad\erratum{a>0}.
\]
Reconstruct the electrical charge distribution that will produce this
potential. Note that $\varphi(\mathbf{r})$ vanishes exponentially for large
$r$, showing that the net charge is zero.
\par\nobreak\hfill\emph{ANS.}\enspace
$\rho(\mathbf{r})=Z\,\delta(\mathbf{r})-\dfrac{Za^{2}e^{-ar}}{4\pi r}$.
\solution
\erratum{The exponential-decay and zero-net-charge conclusion
assumes $a>0$. At $a=0$ the potential is Coulombic and its net charge is
$Z$, not zero.}
Use Poisson's equation in SI units,
\[
\nabla^2\varphi=-\frac{\rho}{\epsilon_0},\qquad
\rho=-\epsilon_0\nabla^2\varphi.
\]
First compute away from the origin. For $h(r)=e^{-ar}/r$,
\[
h'=-e^{-ar}\left(\frac a{r}+\frac1{r^2}\right),\qquad
r^2h'=-(ar+1)e^{-ar},
\]
and therefore
\[
\frac d{dr}(r^2h')
=-a e^{-ar}+a(ar+1)e^{-ar}=a^2r e^{-ar}.
\]
Thus $\nabla^2h=r^{-2}(r^2h')'=a^2e^{-ar}/r$ for $r>0$.
This ordinary calculation omits the point singularity. Its coefficient
is fixed by the flux through a small sphere:
\[
\lim_{R\downarrow0}4\pi R^2h'(R)
=-4\pi\lim_{R\downarrow0}(aR+1)e^{-aR}=-4\pi.
\]
The integral of $a^2h$ over that ball tends to zero, so the full
distributional identity is
\[
\nabla^2\frac{e^{-ar}}r
=a^2\frac{e^{-ar}}r-4\pi\delta^{(3)}(\mathbf r).
\]
Insert $\varphi=Zh/(4\pi\epsilon_0)$ into Poisson's equation:
\[
\boxed{\rho(\mathbf r)=Z\delta^{(3)}(\mathbf r)
-\frac{Za^2e^{-ar}}{4\pi r}.}
\]
The first term is a point charge $Z$ at the origin; the second is a
spherically symmetric cloud of opposite sign. Its $1/r$ density
singularity is integrable because the radial volume element contains
$r^2dr$.

Compute its total charge explicitly. Integration by parts gives
\[
\int_0^\infty r e^{-ar}\,dr
=\left[-\frac r{a}e^{-ar}\right]_0^\infty
+\frac1a\int_0^\infty e^{-ar}\,dr=\frac1{a^2}.
\]
Consequently
\[
Q_{\rm cloud}=4\pi\int_0^\infty r^2
\left(-\frac{Za^2e^{-ar}}{4\pi r}\right)dr=-Z,
\qquad Q_{\rm total}=Z-Z=0.
\]
As an independent check, the enclosed charge at a finite radius is
\[
Q(R)=Z-Za^2\int_0^Rr e^{-ar}\,dr=Z(1+aR)e^{-aR}.
\]
The electric field obtained by differentiating the supplied potential is
\[
E_r=-\varphi'(r)=\frac{Ze^{-ar}}{4\pi\epsilon_0}
\left(\frac a{r}+\frac1{r^2}\right),
\]
and Gauss' law gives $4\pi\epsilon_0R^2E_r(R)=Q(R)$ exactly.
It tends to $Z$ as $R\downarrow0$ and to zero as $R\to\infty$,
confirming both the point charge and the screening cloud.
